\mxfront{terminale-e8}
\mxnewpage\phantomsection\label{sujet:8-1}\addcontentsline{toc}{chapter}{Blanc 2 · Zéro 1}

\section*{Blanc 2 · Zéro 1}

\mxbanner{SUJET ZÉRO · ENTRAÎNEMENT}

\textbf{Durée : 4 h · Note sur 20 · \mxTallied{ex-8-1} exercices indépendants.}\\Calculatrice autorisée. Toute réponse doit être justifiée ; les démarches incomplètes seront prises en compte. Donner les valeurs exactes ; arrondir les probabilités à $10^{-6}$ et les autres approximations à $10^{-4}$, sauf précision différente.

\smallskip\textbf{Nom :}\enspace\dotfill\quad\textbf{Classe :}\enspace\dotfill\quad\textbf{Date :}\enspace\dotfill

\mxtally{ex-8-1}
\subsection*{Exercice 1 · Intégrale et suite \hfill 5 points}Pour tout entier naturel $n$, on pose $I_n=\displaystyle\int_0^1x^ne^{5x}\,\dd x$. Pour $n=0$, l'intégrande est la fonction $x\mapsto e^{5x}$.

\begin{enumerate}

\item Calculer $I_0$. \hfill\textit{(\num{0,5} pt)}

\item À l'aide d'une intégration par parties, calculer $I_1$. \hfill\textit{(\num{0,5} pt)}

\item Démontrer, en comparant les intégrandes, que $(I_n)$ est positive et décroissante. \hfill\textit{(\num{0,5} pt)}

\item Montrer que $I_n\le\frac{e^{5}}{n+1}$ pour tout $n\in\N$. En déduire la limite de $(I_n)$. \hfill\textit{(1 pt)}

\item À l'aide d'une intégration par parties, exprimer $I_{n+1}$ en fonction de $I_n$ et de $n$. \hfill\textit{(\num{0,5} pt)}

\item Montrer que $I_n\ge\frac1{n+1}$ pour tout $n\in\N$. \hfill\textit{(\num{0,5} pt)}

\item On partage $\intff{0}{1}$ en $N$ intervalles de même longueur, avec $N\ge1$. Écrire les sommes $L_N$ et $U_N$ des aires des rectangles construits avec les valeurs de $e^{5x}$ aux extrémités gauches et droites. Justifier $L_N\le I_0\le U_N$, puis calculer $U_N-L_N$. \hfill\textit{(\num{0,5} pt)}

\item Déterminer le plus petit entier $N$ garantissant que cet encadrement a une largeur au plus égale à $10^{-2}$. Écrire une fonction Python qui renvoie $N$, $L_N$ et $U_N$. \hfill\textit{(1 pt)}

\end{enumerate}

\mxtally{ex-8-1}
\subsection*{Exercice 2 · Une fonction logarithmique \hfill 5 points}On considère la fonction $f$ définie sur $\intoo{0}{+\infty}$ par $f(x)=\frac{4+\ln x}x$. On note $\mathcal C_f$ sa courbe représentative.

\begin{enumerate}

\item Déterminer les limites de $f$ en $0^+$ et en $+\infty$. \hfill\textit{(\num{0,5} pt)}

\item Calculer $f' (x)$, puis dresser le tableau de variations de $f$. \hfill\textit{(\num{0,5} pt)}

\item Préciser la valeur du maximum de $f$, puis résoudre $f(x)=0$. \hfill\textit{(\num{0,5} pt)}

\item Étudier la convexité de $f$ et déterminer les coordonnées du point d'inflexion de $\mathcal C_f$. \hfill\textit{(1 pt)}

\item Démontrer que l'équation $f(x)=\frac12 e^{3}$ admet exactement deux solutions dans $\intoo{0}{+\infty}$. \hfill\textit{(\num{0,5} pt)}

\item Déterminer une primitive de $f$ sur $\intoo{0}{+\infty}$. \hfill\textit{(\num{0,5} pt)}

\item Calculer la valeur moyenne de $f$ sur $\intff{1}{e}$. \hfill\textit{(\num{0,5} pt)}

\item Indépendamment des questions précédentes, résoudre $\ln(x-3)+\ln(x+3)=\ln(27)$, après avoir précisé l'ensemble sur lequel cette équation est définie. \hfill\textit{(1 pt)}

\end{enumerate}

\mxnewpage\mxcontinued{Blanc 2 · Zéro 1 · suite}

\mxtally{ex-8-1}
\subsection*{Exercice 3 · Droite, plan et projeté \hfill 5 points}Dans un repère orthonormé de l'espace, on considère le plan $\mathcal P$ d'équation $x+2y+2z=10$ et le point $M\pt{3;0;0}$. La droite $d$ passe par $M$ et admet pour vecteur directeur $\vec n=\pt{1;2;2}$.

\begin{enumerate}

\item Donner une représentation paramétrique de $d$. \hfill\textit{(\num{0,5} pt)}

\item Démontrer que $d$ coupe $\mathcal P$ en un point $H$, puis calculer les coordonnées de $H$. \hfill\textit{(\num{0,5} pt)}

\item Justifier que $H$ est le projeté orthogonal de $M$ sur $\mathcal P$, puis calculer la distance de $M$ à $\mathcal P$. \hfill\textit{(\num{0,5} pt)}

\item On considère le plan $\mathcal Q$ d'équation $y=z$. Démontrer que $\mathcal P\cap\mathcal Q$ est une droite et en donner une représentation paramétrique. \hfill\textit{(1 pt)}

\item Vérifier que $H\in\mathcal Q$. En déduire la distance de $M$ à la droite $\mathcal P\cap\mathcal Q$. \hfill\textit{(\num{0,5} pt)}

\item Déterminer une équation cartésienne du plan parallèle à $\mathcal P$ passant par $M$. \hfill\textit{(\num{0,5} pt)}

\item Déterminer l'angle entre $d$ et la droite $\mathcal P\cap\mathcal Q$. \hfill\textit{(\num{0,5} pt)}

\item On définit le point $N$ par $\vect{MN}=\vec n$. Calculer les distances de $N$ au plan $\mathcal P$ et à la droite $\mathcal P\cap\mathcal Q$, puis les comparer. \hfill\textit{(1 pt)}

\end{enumerate}

\mxtally{ex-8-1}
\subsection*{Exercice 4 · Répétitions et estimation \hfill 5 points}On répète $20$ fois une épreuve à deux issues, appelées succès et échec. Les épreuves sont indépendantes et la probabilité de succès est $p=\frac{3}{10}$ à chaque épreuve. On note $X$ le nombre de succès obtenus.

\begin{enumerate}

\item Justifier que $X$ suit une loi binomiale et préciser ses paramètres. \hfill\textit{(\num{0,5} pt)}

\item Calculer la probabilité d'obtenir exactement deux succès. \hfill\textit{(\num{0,5} pt)}

\item Exprimer à l'aide des événements contraires, puis calculer, la probabilité d'obtenir au moins deux succès. \hfill\textit{(\num{0,5} pt)}

\item Calculer l'espérance et la variance de $X$. \hfill\textit{(1 pt)}

\item On répète maintenant la même épreuve un nombre $N\ge1$ de fois, toujours indépendamment. Déterminer le plus petit entier $N$ pour lequel la probabilité d'obtenir au moins un succès est supérieure ou égale à $\num{0,95}$. Une recherche par tableau ou par boucle est autorisée. \hfill\textit{(\num{0,5} pt)}

\item Pour un entier $j\in\{0,\ldots,20\}$, démontrer la formule de $P(X=j)$ en dénombrant les positions possibles des succès. \hfill\textit{(\num{0,5} pt)}

\item On accepte une série de $20$ épreuves si son nombre de succès est compris entre $1$ et $3$, bornes incluses. Calculer la probabilité qu'une série soit acceptée. \hfill\textit{(\num{0,5} pt)}

\item Écrire une fonction Python sans argument qui renvoie le plus petit entier $j$ tel que $P(X\le j)\ge\num{0,95}$. Exécuter la fonction et vérifier la minimalité du résultat. \hfill\textit{(1 pt)}

\end{enumerate}

\ifmxteacher\mxnewpage

\section*{Corrigé · Blanc 2 · Zéro 1}\label{cor:8-1}

\subsection*{Exercice 1 · Intégrale et suite}\begin{enumerate}

\item \textit{[\num{0,5} pt]} La fonction $x\mapsto \frac{e^{5x}}{5}$ est une primitive de $e^{5x}$. Ainsi $I_0=[\frac{e^{5x}}{5}]_0^1=\frac{e^{5}-1}{5}$.

\item \textit{[\num{0,5} pt]} Choisissons $u(x)=x$ et $v(x)=\frac{e^{5x}}{5}$, deux fonctions dérivables à dérivées continues sur $[0,1]$. Alors $u'=1$ et $v'=e^{5x}$. L'intégration par parties donne
\[\begin{aligned}I_1&=\left[\frac{xe^{5x}}{5}\right]_0^1-\frac{1}{5}\int_0^1e^{5x}\,\dd x\\&=\frac{e^{5}}{5}-\frac{e^{5}-1}{25}=\frac{4e^{5}+1}{25}.\end{aligned}\]

\item \textit{[\num{0,5} pt]} Pour $x\in[0,1]$, $0\le x^{n+1}\le x^n$. Comme $e^{5x}>0$, l'intégration sur des bornes croissantes donne $0\le I_{n+1}\le I_n$. La suite est donc positive et décroissante au sens large. Les inégalités sont même strictes : les intégrandes de $I_n$ et de $I_n-I_{n+1}$ sont continues et strictement positives sur $]0,1[$.

\item \textit{[1 pt]} $e^{5x}\le e^{5}$ ; donc $0\le I_n\le e^{5}\int_0^1x^n\,\dd x=\frac{e^{5}}{n+1}\to0$.

\item \textit{[\num{0,5} pt]} Appliquons l'IPP à $u(x)=x^{n+1}$ et $v(x)=\frac{e^{5x}}{5}$. Comme $n+1\ge1$, le terme en $x=0$ est nul. D'où
\[\begin{aligned}I_{n+1}&=\left[\frac{x^{n+1}e^{5x}}{5}\right]_0^1-\frac{n+1}{5}\int_0^1x^ne^{5x}\,\dd x\\&=\frac{e^{5}}{5}-\frac{n+1}{5}I_n.\end{aligned}\]

\item \textit{[\num{0,5} pt]} Pour $0\le x\le1$, $e^{5x}\ge1$ et $x^n\ge0$. Donc $x^n\le x^ne^{5x}$ ; par intégration,
\[I_n\ge\int_0^1x^n\,\dd x=\left[\frac{x^{n+1}}{n+1}\right]_0^1=\frac{1}{n+1}.\]

\item \textit{[\num{0,5} pt]} La fonction $x\mapsto e^{5x}$ est croissante. Sur chaque intervalle $[\frac{j}{N},\frac{j+1}{N}]$, les valeurs aux extrémités encadrent la fonction. En intégrant puis en sommant,
\[L_N=\frac{1}{N}\sum_{j=0}^{N-1}e^{\frac{5j}{N}}\le I_0\le U_N=\frac{1}{N}\sum_{j=1}^Ne^{\frac{5j}{N}}.\]
Les termes intérieurs s'annulent dans la différence, donc $U_N-L_N=\frac{e^{5}-1}{N}$.

\mxneed{24}

\item \textit{[1 pt]} La largeur est au plus $10^{-2}$ si $N\ge100(e^{5}-1)$. Le premier entier satisfaisant cette borne est $N=14742$. Voici le calcul des bornes :
\begin{lstlisting}
from math import exp, ceil

def bornes():
    N = ceil(100*(exp(5)-1))
    bas = sum(exp(5*j/N) for j in range(N))/N
    haut = bas+(exp(5)-1)/N
    return N, bas, haut
\end{lstlisting}


\end{enumerate}

\subsection*{Exercice 2 · Une fonction logarithmique}\begin{enumerate}

\item \textit{[\num{0,5} pt]} En $0^+$, le numérateur tend vers $-\infty$ et le dénominateur vers $0^+$ : limite $-\infty$ ; en $+\infty$, limite $0$ par croissance comparée.

\item \textit{[\num{0,5} pt]} Par la règle du quotient,
\[f'(x)=\frac{(\frac{1}{x})x-(4+\ln x)}{x^2}=\frac{-3-\ln x}{x^2}.\]
Le dénominateur est positif. Le numérateur s'annule en $s=e^{-3}$, est positif pour $0<x<s$ et négatif pour $x>s$. Ainsi $f$ croît strictement sur $]0,s]$, puis décroît strictement sur $[s,+\infty[$.

\item \textit{[\num{0,5} pt]} Le maximum est atteint uniquement en $s=e^{-3}$ et vaut
\[f(s)=\frac{4- 3}{e^{-3}}=e^{3}.\]
Comme $x>0$, $f(x)=0$ équivaut à $4+\ln x=0$, soit $x=e^{-4}$.

\item \textit{[1 pt]} On dérive $(-3-\ln x)x^{-2}$ :
\[f''(x)=-x^{-3}-2(-3-\ln x)x^{-3}=\frac{2\ln x+5}{x^3}.\]
Le numérateur s'annule en $t=e^{-\frac{5}{2}}$ et passe du signe négatif au signe positif. La fonction est concave sur $]0,t]$, convexe sur $[t,+\infty[$. Elle change de convexité en $t$ et $f(t)=\frac{3}{2t}$ : le point d'inflexion est $I\pt{t;\frac{3}{2t}}$.

\item \textit{[\num{0,5} pt]} Posons $m=\frac{e^{3}}{2}$ et $s=e^{-3}$. La fonction est continue sur $]0,+\infty[$. Sur $]0,s]$, elle est strictement croissante, de limite $-\infty$ en $0$ et de valeur $2m$ en $s$ : le corollaire du TVI donne exactement une solution dans $]0,s[$. Sur $[s,+\infty[$, elle est strictement décroissante de $2m$ vers $0<m$ : exactement une solution dans $]s,+\infty[$. Ces intervalles couvrent le domaine et $f(s)\ne m$ : il y a exactement deux solutions.

\item \textit{[\num{0,5} pt]} Posons $F(x)=4\ln x+\frac12(\ln x)^2$, dérivable pour $x>0$. Alors $F'(x)=\frac{4}{x}+\frac{\ln x}{x}=f(x)$ ; $F$ est une primitive sur $]0,+\infty[$.

\item \textit{[\num{0,5} pt]} Avec cette primitive, $F(e)=\frac{9}{2}$ et $F(1)=0$. L'intégrale vaut donc $\frac{9}{2}$, et la valeur moyenne est
\[\mu=\frac{1}{e-1}\int_1^e f(x)\,\dd x=\frac{9}{2(e-1)}.\]

\item \textit{[1 pt]} Domaine $x>3$ ; $x^2-9=27$, donc $x=6$, la racine négative étant exclue.

\end{enumerate}

\subsection*{Exercice 3 · Droite, plan et projeté}\begin{enumerate}

\item \textit{[\num{0,5} pt]} $\pt{x;y;z}=\pt{3+t;2t;2t}$, $t\in\R$.

\item \textit{[\num{0,5} pt]} $3+9t=10$ donne $t=\frac{7}{9}$ ; $H=\pt{\frac{34}{9};\frac{14}{9};\frac{14}{9}}$.

\item \textit{[\num{0,5} pt]} Le vecteur $\vec n=\pt{1;2;2}$ est normal à $\mathcal P$. Comme $H\in\mathcal P$ et $\vect{MH}=\frac{7}{9}\vec n$, $H$ est le projeté orthogonal de $M$ sur le plan. Le repère étant orthonormé, $\|\vec n\|=\sqrt{1+4+4}=3$. Donc
\[d(M,\mathcal P)=MH=|\frac{7}{9}|\,3=\frac{7}{3}.\]

\item \textit{[1 pt]} Poser $y=z=s$ ; $x=10-4s$. Donc $\pt{x;y;z}=\pt{10;0;0}+s\pt{-4;1;1}$.

\item \textit{[\num{0,5} pt]} Oui, ses deux dernières coordonnées sont égales. $H$ est dans la droite et $\vect{MH}$ est orthogonal à toute direction du plan, donc à cette droite. Distance $\frac{7}{3}$.

\item \textit{[\num{0,5} pt]} $x+2y+2z=3$ ; même vecteur normal et passage par $M$.

\item \textit{[\num{0,5} pt]} $\vec n\cdot\pt{-4;1;1}=-4+2+2=0$ : angle $\frac{\pi}{2}$.

\item \textit{[1 pt]} $N=M+\vec n$ est sur la même droite normale que $M$ ; son projeté sur $\mathcal P$ est donc aussi $H$. Or $H\in\mathcal P\cap\mathcal Q$. Pour tout $P$ de cette droite, $NP^2=NH^2+HP^2$, donc sa distance minimale à $N$ est $NH$. Finalement,
\[d(N,\mathcal P)=d(N,\mathcal P\cap\mathcal Q)=\|\vect{NH}\|=|\frac{7}{9}-1|\,\|\vec n\|=\frac{2}{3}.\]

\end{enumerate}

\subsection*{Exercice 4 · Répétitions et estimation}\begin{enumerate}

\item \textit{[\num{0,5} pt]} Répétition d'un nombre fixe $20$ d'épreuves identiques indépendantes à deux issues ; $X\sim\mathcal B\pt{20;\frac{3}{10}}$.

\item \textit{[\num{0,5} pt]} $P(X=2)=\binom{20}2(\frac{3}{10})^2(\frac{7}{10})^{18}\simeq\num{0,027846}$.

\item \textit{[\num{0,5} pt]} $1-P(X=0)-P(X=1)=1-(\frac{7}{10})^{20}-20(\frac{3}{10})(\frac{7}{10})^{19}\simeq\num{0,992363}$.

\item \textit{[1 pt]} $E(X)=6$, $V(X)=\frac{21}{5}$.

\item \textit{[\num{0,5} pt]} Pour $N$ essais indépendants, la probabilité d'aucun succès est $(1-p)^N$. La probabilité demandée est donc $1-(\frac{7}{10})^N$. Elle est strictement croissante avec $N$, car $0<\frac{7}{10}<1$. Les calculs donnent
\[P_{8}\simeq\num{0,942352}<\num{0,95},\qquad P_{9}\simeq\num{0,959646}\ge\num{0,95}.\]
Ainsi $9$ est le plus petit nombre d'essais convenable.

\item \textit{[\num{0,5} pt]} Notons $n=20$. Pour $j\in\{0,\ldots,n\}$, un résultat ordonné contenant $j$ succès et $n-j$ échecs a, par indépendance, la probabilité $p^j(1-p)^{n-j}$. Il existe $\binom nj$ choix des positions des succès. Ces résultats sont incompatibles, donc
\[P(X=j)=\binom nj p^j(1-p)^{n-j},\qquad p=\frac{3}{10}.\]

\item \textit{[\num{0,5} pt]} Les événements $X=1$, $X=2$, $X=3$ sont incompatibles, donc
\[P(1\le X\le3)=\sum_{j=1}^3\binom{20}j(\frac{3}{10})^j(\frac{7}{10})^{20-j}\simeq\num{0,106289}.\]

\mxneed{24}

\item \textit{[1 pt]} 
\begin{lstlisting}
from math import comb

def quantile():
    n, p = 20, 3/10
    j, cumul = 0, (1-p)**n
    while cumul < 0.95:
        j += 1
        cumul += comb(n, j)*p**j*(1-p)**(n-j)
    return j
\end{lstlisting}
Les probabilités ajoutées sont positives et le cumul final vaut $1$, donc la boucle s'arrête. Elle renvoie $9$. Les cumuls aux rangs $8$ et $9$ sont environ $\num{0,886669}$ et $\num{0,952038}$ : cela confirme la minimalité.

\end{enumerate}

\fi

\mxnewpage\phantomsection\label{sujet:8-2}\addcontentsline{toc}{chapter}{Blanc 2 · Zéro 2}

\section*{Blanc 2 · Zéro 2}

\mxbanner{SUJET ZÉRO · ENTRAÎNEMENT}

\textbf{Durée : 4 h · Note sur 20 · \mxTallied{ex-8-2} exercices indépendants.}\\Calculatrice autorisée. Toute réponse doit être justifiée ; les démarches incomplètes seront prises en compte. Donner les valeurs exactes ; arrondir les probabilités à $10^{-6}$ et les autres approximations à $10^{-4}$, sauf précision différente.

\smallskip\textbf{Nom :}\enspace\dotfill\quad\textbf{Classe :}\enspace\dotfill\quad\textbf{Date :}\enspace\dotfill

\mxtally{ex-8-2}
\subsection*{Exercice 1 · Intégrale et suite \hfill 5 points}Pour tout entier naturel $n$, on pose $I_n=\displaystyle\int_0^1x^ne^{4x}\,\dd x$. Pour $n=0$, l'intégrande est la fonction $x\mapsto e^{4x}$.

\begin{enumerate}

\item Calculer $I_0$. \hfill\textit{(\num{0,5} pt)}

\item À l'aide d'une intégration par parties, calculer $I_1$. \hfill\textit{(\num{0,5} pt)}

\item Démontrer, en comparant les intégrandes, que $(I_n)$ est positive et décroissante. \hfill\textit{(\num{0,5} pt)}

\item Montrer que $I_n\le\frac{e^{4}}{n+1}$ pour tout $n\in\N$. En déduire la limite de $(I_n)$. \hfill\textit{(1 pt)}

\item À l'aide d'une intégration par parties, exprimer $I_{n+1}$ en fonction de $I_n$ et de $n$. \hfill\textit{(\num{0,5} pt)}

\item Montrer que $I_n\ge\frac1{n+1}$ pour tout $n\in\N$. \hfill\textit{(\num{0,5} pt)}

\item On partage $\intff{0}{1}$ en $N$ intervalles de même longueur, avec $N\ge1$. Écrire les sommes $L_N$ et $U_N$ des aires des rectangles construits avec les valeurs de $e^{4x}$ aux extrémités gauches et droites. Justifier $L_N\le I_0\le U_N$, puis calculer $U_N-L_N$. \hfill\textit{(\num{0,5} pt)}

\item Déterminer le plus petit entier $N$ garantissant que cet encadrement a une largeur au plus égale à $10^{-2}$. Écrire une fonction Python qui renvoie $N$, $L_N$ et $U_N$. \hfill\textit{(1 pt)}

\end{enumerate}

\mxtally{ex-8-2}
\subsection*{Exercice 2 · Une fonction logarithmique \hfill 5 points}On considère la fonction $f$ définie sur $\intoo{0}{+\infty}$ par $f(x)=\frac{3+\ln x}x$. On note $\mathcal C_f$ sa courbe représentative.

\begin{enumerate}

\item Déterminer les limites de $f$ en $0^+$ et en $+\infty$. \hfill\textit{(\num{0,5} pt)}

\item Calculer $f' (x)$, puis dresser le tableau de variations de $f$. \hfill\textit{(\num{0,5} pt)}

\item Préciser la valeur du maximum de $f$, puis résoudre $f(x)=0$. \hfill\textit{(\num{0,5} pt)}

\item Étudier la convexité de $f$ et déterminer les coordonnées du point d'inflexion de $\mathcal C_f$. \hfill\textit{(1 pt)}

\item Démontrer que l'équation $f(x)=\frac12 e^{2}$ admet exactement deux solutions dans $\intoo{0}{+\infty}$. \hfill\textit{(\num{0,5} pt)}

\item Déterminer une primitive de $f$ sur $\intoo{0}{+\infty}$. \hfill\textit{(\num{0,5} pt)}

\item Calculer la valeur moyenne de $f$ sur $\intff{1}{e}$. \hfill\textit{(\num{0,5} pt)}

\item Indépendamment des questions précédentes, résoudre $\ln(x-1)+\ln(x+1)=\ln(63)$, après avoir précisé l'ensemble sur lequel cette équation est définie. \hfill\textit{(1 pt)}

\end{enumerate}

\mxnewpage\mxcontinued{Blanc 2 · Zéro 2 · suite}

\mxtally{ex-8-2}
\subsection*{Exercice 3 · Droite, plan et projeté \hfill 5 points}Dans un repère orthonormé de l'espace, on considère le plan $\mathcal P$ d'équation $x+2y+2z=6$ et le point $M\pt{1;0;0}$. La droite $d$ passe par $M$ et admet pour vecteur directeur $\vec n=\pt{1;2;2}$.

\begin{enumerate}

\item Donner une représentation paramétrique de $d$. \hfill\textit{(\num{0,5} pt)}

\item Démontrer que $d$ coupe $\mathcal P$ en un point $H$, puis calculer les coordonnées de $H$. \hfill\textit{(\num{0,5} pt)}

\item Justifier que $H$ est le projeté orthogonal de $M$ sur $\mathcal P$, puis calculer la distance de $M$ à $\mathcal P$. \hfill\textit{(\num{0,5} pt)}

\item On considère le plan $\mathcal Q$ d'équation $y=z$. Démontrer que $\mathcal P\cap\mathcal Q$ est une droite et en donner une représentation paramétrique. \hfill\textit{(1 pt)}

\item Vérifier que $H\in\mathcal Q$. En déduire la distance de $M$ à la droite $\mathcal P\cap\mathcal Q$. \hfill\textit{(\num{0,5} pt)}

\item Déterminer une équation cartésienne du plan parallèle à $\mathcal P$ passant par $M$. \hfill\textit{(\num{0,5} pt)}

\item Déterminer l'angle entre $d$ et la droite $\mathcal P\cap\mathcal Q$. \hfill\textit{(\num{0,5} pt)}

\item On définit le point $N$ par $\vect{MN}=\vec n$. Calculer les distances de $N$ au plan $\mathcal P$ et à la droite $\mathcal P\cap\mathcal Q$, puis les comparer. \hfill\textit{(1 pt)}

\end{enumerate}

\mxtally{ex-8-2}
\subsection*{Exercice 4 · Répétitions et estimation \hfill 5 points}On répète $17$ fois une épreuve à deux issues, appelées succès et échec. Les épreuves sont indépendantes et la probabilité de succès est $p=\frac{3}{10}$ à chaque épreuve. On note $X$ le nombre de succès obtenus.

\begin{enumerate}

\item Justifier que $X$ suit une loi binomiale et préciser ses paramètres. \hfill\textit{(\num{0,5} pt)}

\item Calculer la probabilité d'obtenir exactement deux succès. \hfill\textit{(\num{0,5} pt)}

\item Exprimer à l'aide des événements contraires, puis calculer, la probabilité d'obtenir au moins deux succès. \hfill\textit{(\num{0,5} pt)}

\item Calculer l'espérance et la variance de $X$. \hfill\textit{(1 pt)}

\item On répète maintenant la même épreuve un nombre $N\ge1$ de fois, toujours indépendamment. Déterminer le plus petit entier $N$ pour lequel la probabilité d'obtenir au moins un succès est supérieure ou égale à $\num{0,95}$. Une recherche par tableau ou par boucle est autorisée. \hfill\textit{(\num{0,5} pt)}

\item Pour un entier $j\in\{0,\ldots,17\}$, démontrer la formule de $P(X=j)$ en dénombrant les positions possibles des succès. \hfill\textit{(\num{0,5} pt)}

\item On accepte une série de $17$ épreuves si son nombre de succès est compris entre $1$ et $3$, bornes incluses. Calculer la probabilité qu'une série soit acceptée. \hfill\textit{(\num{0,5} pt)}

\item Écrire une fonction Python sans argument qui renvoie le plus petit entier $j$ tel que $P(X\le j)\ge\num{0,95}$. Exécuter la fonction et vérifier la minimalité du résultat. \hfill\textit{(1 pt)}

\end{enumerate}

\ifmxteacher\mxnewpage

\section*{Corrigé · Blanc 2 · Zéro 2}\label{cor:8-2}

\subsection*{Exercice 1 · Intégrale et suite}\begin{enumerate}

\item \textit{[\num{0,5} pt]} La fonction $x\mapsto \frac{e^{4x}}{4}$ est une primitive de $e^{4x}$. Ainsi $I_0=[\frac{e^{4x}}{4}]_0^1=\frac{e^{4}-1}{4}$.

\item \textit{[\num{0,5} pt]} Choisissons $u(x)=x$ et $v(x)=\frac{e^{4x}}{4}$, deux fonctions dérivables à dérivées continues sur $[0,1]$. Alors $u'=1$ et $v'=e^{4x}$. L'intégration par parties donne
\[\begin{aligned}I_1&=\left[\frac{xe^{4x}}{4}\right]_0^1-\frac{1}{4}\int_0^1e^{4x}\,\dd x\\&=\frac{e^{4}}{4}-\frac{e^{4}-1}{16}=\frac{3e^{4}+1}{16}.\end{aligned}\]

\item \textit{[\num{0,5} pt]} Pour $x\in[0,1]$, $0\le x^{n+1}\le x^n$. Comme $e^{4x}>0$, l'intégration sur des bornes croissantes donne $0\le I_{n+1}\le I_n$. La suite est donc positive et décroissante au sens large. Les inégalités sont même strictes : les intégrandes de $I_n$ et de $I_n-I_{n+1}$ sont continues et strictement positives sur $]0,1[$.

\item \textit{[1 pt]} $e^{4x}\le e^{4}$ ; donc $0\le I_n\le e^{4}\int_0^1x^n\,\dd x=\frac{e^{4}}{n+1}\to0$.

\item \textit{[\num{0,5} pt]} Appliquons l'IPP à $u(x)=x^{n+1}$ et $v(x)=\frac{e^{4x}}{4}$. Comme $n+1\ge1$, le terme en $x=0$ est nul. D'où
\[\begin{aligned}I_{n+1}&=\left[\frac{x^{n+1}e^{4x}}{4}\right]_0^1-\frac{n+1}{4}\int_0^1x^ne^{4x}\,\dd x\\&=\frac{e^{4}}{4}-\frac{n+1}{4}I_n.\end{aligned}\]

\item \textit{[\num{0,5} pt]} Pour $0\le x\le1$, $e^{4x}\ge1$ et $x^n\ge0$. Donc $x^n\le x^ne^{4x}$ ; par intégration,
\[I_n\ge\int_0^1x^n\,\dd x=\left[\frac{x^{n+1}}{n+1}\right]_0^1=\frac{1}{n+1}.\]

\item \textit{[\num{0,5} pt]} La fonction $x\mapsto e^{4x}$ est croissante. Sur chaque intervalle $[\frac{j}{N},\frac{j+1}{N}]$, les valeurs aux extrémités encadrent la fonction. En intégrant puis en sommant,
\[L_N=\frac{1}{N}\sum_{j=0}^{N-1}e^{\frac{4j}{N}}\le I_0\le U_N=\frac{1}{N}\sum_{j=1}^Ne^{\frac{4j}{N}}.\]
Les termes intérieurs s'annulent dans la différence, donc $U_N-L_N=\frac{e^{4}-1}{N}$.

\mxneed{24}

\item \textit{[1 pt]} La largeur est au plus $10^{-2}$ si $N\ge100(e^{4}-1)$. Le premier entier satisfaisant cette borne est $N=5360$. Voici le calcul des bornes :
\begin{lstlisting}
from math import exp, ceil

def bornes():
    N = ceil(100*(exp(4)-1))
    bas = sum(exp(4*j/N) for j in range(N))/N
    haut = bas+(exp(4)-1)/N
    return N, bas, haut
\end{lstlisting}


\end{enumerate}

\subsection*{Exercice 2 · Une fonction logarithmique}\begin{enumerate}

\item \textit{[\num{0,5} pt]} En $0^+$, le numérateur tend vers $-\infty$ et le dénominateur vers $0^+$ : limite $-\infty$ ; en $+\infty$, limite $0$ par croissance comparée.

\item \textit{[\num{0,5} pt]} Par la règle du quotient,
\[f'(x)=\frac{(\frac{1}{x})x-(3+\ln x)}{x^2}=\frac{-2-\ln x}{x^2}.\]
Le dénominateur est positif. Le numérateur s'annule en $s=e^{-2}$, est positif pour $0<x<s$ et négatif pour $x>s$. Ainsi $f$ croît strictement sur $]0,s]$, puis décroît strictement sur $[s,+\infty[$.

\item \textit{[\num{0,5} pt]} Le maximum est atteint uniquement en $s=e^{-2}$ et vaut
\[f(s)=\frac{3- 2}{e^{-2}}=e^{2}.\]
Comme $x>0$, $f(x)=0$ équivaut à $3+\ln x=0$, soit $x=e^{-3}$.

\item \textit{[1 pt]} On dérive $(-2-\ln x)x^{-2}$ :
\[f''(x)=-x^{-3}-2(-2-\ln x)x^{-3}=\frac{2\ln x+3}{x^3}.\]
Le numérateur s'annule en $t=e^{-\frac{3}{2}}$ et passe du signe négatif au signe positif. La fonction est concave sur $]0,t]$, convexe sur $[t,+\infty[$. Elle change de convexité en $t$ et $f(t)=\frac{3}{2t}$ : le point d'inflexion est $I\pt{t;\frac{3}{2t}}$.

\item \textit{[\num{0,5} pt]} Posons $m=\frac{e^{2}}{2}$ et $s=e^{-2}$. La fonction est continue sur $]0,+\infty[$. Sur $]0,s]$, elle est strictement croissante, de limite $-\infty$ en $0$ et de valeur $2m$ en $s$ : le corollaire du TVI donne exactement une solution dans $]0,s[$. Sur $[s,+\infty[$, elle est strictement décroissante de $2m$ vers $0<m$ : exactement une solution dans $]s,+\infty[$. Ces intervalles couvrent le domaine et $f(s)\ne m$ : il y a exactement deux solutions.

\item \textit{[\num{0,5} pt]} Posons $F(x)=3\ln x+\frac12(\ln x)^2$, dérivable pour $x>0$. Alors $F'(x)=\frac{3}{x}+\frac{\ln x}{x}=f(x)$ ; $F$ est une primitive sur $]0,+\infty[$.

\item \textit{[\num{0,5} pt]} Avec cette primitive, $F(e)=\frac{7}{2}$ et $F(1)=0$. L'intégrale vaut donc $\frac{7}{2}$, et la valeur moyenne est
\[\mu=\frac{1}{e-1}\int_1^e f(x)\,\dd x=\frac{7}{2(e-1)}.\]

\item \textit{[1 pt]} Domaine $x>1$ ; $x^2-1=63$, donc $x=8$, la racine négative étant exclue.

\end{enumerate}

\subsection*{Exercice 3 · Droite, plan et projeté}\begin{enumerate}

\item \textit{[\num{0,5} pt]} $\pt{x;y;z}=\pt{1+t;2t;2t}$, $t\in\R$.

\item \textit{[\num{0,5} pt]} $1+9t=6$ donne $t=\frac{5}{9}$ ; $H=\pt{\frac{14}{9};\frac{10}{9};\frac{10}{9}}$.

\item \textit{[\num{0,5} pt]} Le vecteur $\vec n=\pt{1;2;2}$ est normal à $\mathcal P$. Comme $H\in\mathcal P$ et $\vect{MH}=\frac{5}{9}\vec n$, $H$ est le projeté orthogonal de $M$ sur le plan. Le repère étant orthonormé, $\|\vec n\|=\sqrt{1+4+4}=3$. Donc
\[d(M,\mathcal P)=MH=|\frac{5}{9}|\,3=\frac{5}{3}.\]

\item \textit{[1 pt]} Poser $y=z=s$ ; $x=6-4s$. Donc $\pt{x;y;z}=\pt{6;0;0}+s\pt{-4;1;1}$.

\item \textit{[\num{0,5} pt]} Oui, ses deux dernières coordonnées sont égales. $H$ est dans la droite et $\vect{MH}$ est orthogonal à toute direction du plan, donc à cette droite. Distance $\frac{5}{3}$.

\item \textit{[\num{0,5} pt]} $x+2y+2z=1$ ; même vecteur normal et passage par $M$.

\item \textit{[\num{0,5} pt]} $\vec n\cdot\pt{-4;1;1}=-4+2+2=0$ : angle $\frac{\pi}{2}$.

\item \textit{[1 pt]} $N=M+\vec n$ est sur la même droite normale que $M$ ; son projeté sur $\mathcal P$ est donc aussi $H$. Or $H\in\mathcal P\cap\mathcal Q$. Pour tout $P$ de cette droite, $NP^2=NH^2+HP^2$, donc sa distance minimale à $N$ est $NH$. Finalement,
\[d(N,\mathcal P)=d(N,\mathcal P\cap\mathcal Q)=\|\vect{NH}\|=|\frac{5}{9}-1|\,\|\vec n\|=\frac{4}{3}.\]

\end{enumerate}

\subsection*{Exercice 4 · Répétitions et estimation}\begin{enumerate}

\item \textit{[\num{0,5} pt]} Répétition d'un nombre fixe $17$ d'épreuves identiques indépendantes à deux issues ; $X\sim\mathcal B\pt{17;\frac{3}{10}}$.

\item \textit{[\num{0,5} pt]} $P(X=2)=\binom{17}2(\frac{3}{10})^2(\frac{7}{10})^{15}\simeq\num{0,058110}$.

\item \textit{[\num{0,5} pt]} $1-P(X=0)-P(X=1)=1-(\frac{7}{10})^{17}-17(\frac{3}{10})(\frac{7}{10})^{16}\simeq\num{0,980725}$.

\item \textit{[1 pt]} $E(X)=\frac{51}{10}$, $V(X)=\frac{357}{100}$.

\item \textit{[\num{0,5} pt]} Pour $N$ essais indépendants, la probabilité d'aucun succès est $(1-p)^N$. La probabilité demandée est donc $1-(\frac{7}{10})^N$. Elle est strictement croissante avec $N$, car $0<\frac{7}{10}<1$. Les calculs donnent
\[P_{8}\simeq\num{0,942352}<\num{0,95},\qquad P_{9}\simeq\num{0,959646}\ge\num{0,95}.\]
Ainsi $9$ est le plus petit nombre d'essais convenable.

\item \textit{[\num{0,5} pt]} Notons $n=17$. Pour $j\in\{0,\ldots,n\}$, un résultat ordonné contenant $j$ succès et $n-j$ échecs a, par indépendance, la probabilité $p^j(1-p)^{n-j}$. Il existe $\binom nj$ choix des positions des succès. Ces résultats sont incompatibles, donc
\[P(X=j)=\binom nj p^j(1-p)^{n-j},\qquad p=\frac{3}{10}.\]

\item \textit{[\num{0,5} pt]} Les événements $X=1$, $X=2$, $X=3$ sont incompatibles, donc
\[P(1\le X\le3)=\sum_{j=1}^3\binom{17}j(\frac{3}{10})^j(\frac{7}{10})^{17-j}\simeq\num{0,199581}.\]

\mxneed{24}

\item \textit{[1 pt]} 
\begin{lstlisting}
from math import comb

def quantile():
    n, p = 17, 3/10
    j, cumul = 0, (1-p)**n
    while cumul < 0.95:
        j += 1
        cumul += comb(n, j)*p**j*(1-p)**(n-j)
    return j
\end{lstlisting}
Les probabilités ajoutées sont positives et le cumul final vaut $1$, donc la boucle s'arrête. Elle renvoie $8$. Les cumuls aux rangs $7$ et $8$ sont environ $\num{0,895360}$ et $\num{0,959723}$ : cela confirme la minimalité.

\end{enumerate}

\fi

\ifmxteacher

\mxnewpage\phantomsection\label{sujet:8-3}\addcontentsline{toc}{chapter}{Blanc 2 · A}

\section*{Blanc 2 · A}

\mxbanner*{SUJET RÉSERVÉ À L'ÉVALUATION}

\textbf{Durée : 4 h · Note sur 20 · \mxTallied{ex-8-3} exercices indépendants.}\\Calculatrice autorisée. Toute réponse doit être justifiée ; les démarches incomplètes seront prises en compte. Donner les valeurs exactes ; arrondir les probabilités à $10^{-6}$ et les autres approximations à $10^{-4}$, sauf précision différente.

\smallskip\textbf{Nom :}\enspace\dotfill\quad\textbf{Classe :}\enspace\dotfill\quad\textbf{Date :}\enspace\dotfill

\mxtally{ex-8-3}
\subsection*{Exercice 1 · Intégrale et suite \hfill 5 points}Pour tout entier naturel $n$, on pose $I_n=\displaystyle\int_0^1x^ne^{3x}\,\dd x$. Pour $n=0$, l'intégrande est la fonction $x\mapsto e^{3x}$.

\begin{enumerate}

\item Calculer $I_0$. \hfill\textit{(\num{0,5} pt)}

\item À l'aide d'une intégration par parties, calculer $I_1$. \hfill\textit{(\num{0,5} pt)}

\item Démontrer, en comparant les intégrandes, que $(I_n)$ est positive et décroissante. \hfill\textit{(\num{0,5} pt)}

\item Montrer que $I_n\le\frac{e^{3}}{n+1}$ pour tout $n\in\N$. En déduire la limite de $(I_n)$. \hfill\textit{(1 pt)}

\item À l'aide d'une intégration par parties, exprimer $I_{n+1}$ en fonction de $I_n$ et de $n$. \hfill\textit{(\num{0,5} pt)}

\item Montrer que $I_n\ge\frac1{n+1}$ pour tout $n\in\N$. \hfill\textit{(\num{0,5} pt)}

\item On partage $\intff{0}{1}$ en $N$ intervalles de même longueur, avec $N\ge1$. Écrire les sommes $L_N$ et $U_N$ des aires des rectangles construits avec les valeurs de $e^{3x}$ aux extrémités gauches et droites. Justifier $L_N\le I_0\le U_N$, puis calculer $U_N-L_N$. \hfill\textit{(\num{0,5} pt)}

\item Déterminer le plus petit entier $N$ garantissant que cet encadrement a une largeur au plus égale à $10^{-2}$. Écrire une fonction Python qui renvoie $N$, $L_N$ et $U_N$. \hfill\textit{(1 pt)}

\end{enumerate}

\mxtally{ex-8-3}
\subsection*{Exercice 2 · Une fonction logarithmique \hfill 5 points}On considère la fonction $f$ définie sur $\intoo{0}{+\infty}$ par $f(x)=\frac{2+\ln x}x$. On note $\mathcal C_f$ sa courbe représentative.

\begin{enumerate}

\item Déterminer les limites de $f$ en $0^+$ et en $+\infty$. \hfill\textit{(\num{0,5} pt)}

\item Calculer $f' (x)$, puis dresser le tableau de variations de $f$. \hfill\textit{(\num{0,5} pt)}

\item Préciser la valeur du maximum de $f$, puis résoudre $f(x)=0$. \hfill\textit{(\num{0,5} pt)}

\item Étudier la convexité de $f$ et déterminer les coordonnées du point d'inflexion de $\mathcal C_f$. \hfill\textit{(1 pt)}

\item Démontrer que l'équation $f(x)=\frac12 e^{1}$ admet exactement deux solutions dans $\intoo{0}{+\infty}$. \hfill\textit{(\num{0,5} pt)}

\item Déterminer une primitive de $f$ sur $\intoo{0}{+\infty}$. \hfill\textit{(\num{0,5} pt)}

\item Calculer la valeur moyenne de $f$ sur $\intff{1}{e}$. \hfill\textit{(\num{0,5} pt)}

\item Indépendamment des questions précédentes, résoudre $\ln(x-2)+\ln(x+2)=\ln(21)$, après avoir précisé l'ensemble sur lequel cette équation est définie. \hfill\textit{(1 pt)}

\end{enumerate}

\mxnewpage\mxcontinued{Blanc 2 · A · suite}

\mxtally{ex-8-3}
\subsection*{Exercice 3 · Droite, plan et projeté \hfill 5 points}Dans un repère orthonormé de l'espace, on considère le plan $\mathcal P$ d'équation $x+2y+2z=10$ et le point $M\pt{2;0;0}$. La droite $d$ passe par $M$ et admet pour vecteur directeur $\vec n=\pt{1;2;2}$.

\begin{enumerate}

\item Donner une représentation paramétrique de $d$. \hfill\textit{(\num{0,5} pt)}

\item Démontrer que $d$ coupe $\mathcal P$ en un point $H$, puis calculer les coordonnées de $H$. \hfill\textit{(\num{0,5} pt)}

\item Justifier que $H$ est le projeté orthogonal de $M$ sur $\mathcal P$, puis calculer la distance de $M$ à $\mathcal P$. \hfill\textit{(\num{0,5} pt)}

\item On considère le plan $\mathcal Q$ d'équation $y=z$. Démontrer que $\mathcal P\cap\mathcal Q$ est une droite et en donner une représentation paramétrique. \hfill\textit{(1 pt)}

\item Vérifier que $H\in\mathcal Q$. En déduire la distance de $M$ à la droite $\mathcal P\cap\mathcal Q$. \hfill\textit{(\num{0,5} pt)}

\item Déterminer une équation cartésienne du plan parallèle à $\mathcal P$ passant par $M$. \hfill\textit{(\num{0,5} pt)}

\item Déterminer l'angle entre $d$ et la droite $\mathcal P\cap\mathcal Q$. \hfill\textit{(\num{0,5} pt)}

\item On définit le point $N$ par $\vect{MN}=\vec n$. Calculer les distances de $N$ au plan $\mathcal P$ et à la droite $\mathcal P\cap\mathcal Q$, puis les comparer. \hfill\textit{(1 pt)}

\end{enumerate}

\mxtally{ex-8-3}
\subsection*{Exercice 4 · Répétitions et estimation \hfill 5 points}On répète $18$ fois une épreuve à deux issues, appelées succès et échec. Les épreuves sont indépendantes et la probabilité de succès est $p=\frac{1}{5}$ à chaque épreuve. On note $X$ le nombre de succès obtenus.

\begin{enumerate}

\item Justifier que $X$ suit une loi binomiale et préciser ses paramètres. \hfill\textit{(\num{0,5} pt)}

\item Calculer la probabilité d'obtenir exactement deux succès. \hfill\textit{(\num{0,5} pt)}

\item Exprimer à l'aide des événements contraires, puis calculer, la probabilité d'obtenir au moins deux succès. \hfill\textit{(\num{0,5} pt)}

\item Calculer l'espérance et la variance de $X$. \hfill\textit{(1 pt)}

\item On répète maintenant la même épreuve un nombre $N\ge1$ de fois, toujours indépendamment. Déterminer le plus petit entier $N$ pour lequel la probabilité d'obtenir au moins un succès est supérieure ou égale à $\num{0,95}$. Une recherche par tableau ou par boucle est autorisée. \hfill\textit{(\num{0,5} pt)}

\item Pour un entier $j\in\{0,\ldots,18\}$, démontrer la formule de $P(X=j)$ en dénombrant les positions possibles des succès. \hfill\textit{(\num{0,5} pt)}

\item On accepte une série de $18$ épreuves si son nombre de succès est compris entre $1$ et $3$, bornes incluses. Calculer la probabilité qu'une série soit acceptée. \hfill\textit{(\num{0,5} pt)}

\item Écrire une fonction Python sans argument qui renvoie le plus petit entier $j$ tel que $P(X\le j)\ge\num{0,95}$. Exécuter la fonction et vérifier la minimalité du résultat. \hfill\textit{(1 pt)}

\end{enumerate}

\ifmxteacher\mxnewpage

\section*{Corrigé · Blanc 2 · A}\label{cor:8-3}

\subsection*{Exercice 1 · Intégrale et suite}\begin{enumerate}

\item \textit{[\num{0,5} pt]} La fonction $x\mapsto \frac{e^{3x}}{3}$ est une primitive de $e^{3x}$. Ainsi $I_0=[\frac{e^{3x}}{3}]_0^1=\frac{e^{3}-1}{3}$.

\item \textit{[\num{0,5} pt]} Choisissons $u(x)=x$ et $v(x)=\frac{e^{3x}}{3}$, deux fonctions dérivables à dérivées continues sur $[0,1]$. Alors $u'=1$ et $v'=e^{3x}$. L'intégration par parties donne
\[\begin{aligned}I_1&=\left[\frac{xe^{3x}}{3}\right]_0^1-\frac{1}{3}\int_0^1e^{3x}\,\dd x\\&=\frac{e^{3}}{3}-\frac{e^{3}-1}{9}=\frac{2e^{3}+1}{9}.\end{aligned}\]

\item \textit{[\num{0,5} pt]} Pour $x\in[0,1]$, $0\le x^{n+1}\le x^n$. Comme $e^{3x}>0$, l'intégration sur des bornes croissantes donne $0\le I_{n+1}\le I_n$. La suite est donc positive et décroissante au sens large. Les inégalités sont même strictes : les intégrandes de $I_n$ et de $I_n-I_{n+1}$ sont continues et strictement positives sur $]0,1[$.

\item \textit{[1 pt]} $e^{3x}\le e^{3}$ ; donc $0\le I_n\le e^{3}\int_0^1x^n\,\dd x=\frac{e^{3}}{n+1}\to0$.

\item \textit{[\num{0,5} pt]} Appliquons l'IPP à $u(x)=x^{n+1}$ et $v(x)=\frac{e^{3x}}{3}$. Comme $n+1\ge1$, le terme en $x=0$ est nul. D'où
\[\begin{aligned}I_{n+1}&=\left[\frac{x^{n+1}e^{3x}}{3}\right]_0^1-\frac{n+1}{3}\int_0^1x^ne^{3x}\,\dd x\\&=\frac{e^{3}}{3}-\frac{n+1}{3}I_n.\end{aligned}\]

\item \textit{[\num{0,5} pt]} Pour $0\le x\le1$, $e^{3x}\ge1$ et $x^n\ge0$. Donc $x^n\le x^ne^{3x}$ ; par intégration,
\[I_n\ge\int_0^1x^n\,\dd x=\left[\frac{x^{n+1}}{n+1}\right]_0^1=\frac{1}{n+1}.\]

\item \textit{[\num{0,5} pt]} La fonction $x\mapsto e^{3x}$ est croissante. Sur chaque intervalle $[\frac{j}{N},\frac{j+1}{N}]$, les valeurs aux extrémités encadrent la fonction. En intégrant puis en sommant,
\[L_N=\frac{1}{N}\sum_{j=0}^{N-1}e^{\frac{3j}{N}}\le I_0\le U_N=\frac{1}{N}\sum_{j=1}^Ne^{\frac{3j}{N}}.\]
Les termes intérieurs s'annulent dans la différence, donc $U_N-L_N=\frac{e^{3}-1}{N}$.

\mxneed{24}

\item \textit{[1 pt]} La largeur est au plus $10^{-2}$ si $N\ge100(e^{3}-1)$. Le premier entier satisfaisant cette borne est $N=1909$. Voici le calcul des bornes :
\begin{lstlisting}
from math import exp, ceil

def bornes():
    N = ceil(100*(exp(3)-1))
    bas = sum(exp(3*j/N) for j in range(N))/N
    haut = bas+(exp(3)-1)/N
    return N, bas, haut
\end{lstlisting}


\end{enumerate}

\subsection*{Exercice 2 · Une fonction logarithmique}\begin{enumerate}

\item \textit{[\num{0,5} pt]} En $0^+$, le numérateur tend vers $-\infty$ et le dénominateur vers $0^+$ : limite $-\infty$ ; en $+\infty$, limite $0$ par croissance comparée.

\item \textit{[\num{0,5} pt]} Par la règle du quotient,
\[f'(x)=\frac{(\frac{1}{x})x-(2+\ln x)}{x^2}=\frac{-1-\ln x}{x^2}.\]
Le dénominateur est positif. Le numérateur s'annule en $s=e^{-1}$, est positif pour $0<x<s$ et négatif pour $x>s$. Ainsi $f$ croît strictement sur $]0,s]$, puis décroît strictement sur $[s,+\infty[$.

\item \textit{[\num{0,5} pt]} Le maximum est atteint uniquement en $s=e^{-1}$ et vaut
\[f(s)=\frac{2- 1}{e^{-1}}=e^{1}.\]
Comme $x>0$, $f(x)=0$ équivaut à $2+\ln x=0$, soit $x=e^{-2}$.

\item \textit{[1 pt]} On dérive $(-1-\ln x)x^{-2}$ :
\[f''(x)=-x^{-3}-2(-1-\ln x)x^{-3}=\frac{2\ln x+1}{x^3}.\]
Le numérateur s'annule en $t=e^{-\frac{1}{2}}$ et passe du signe négatif au signe positif. La fonction est concave sur $]0,t]$, convexe sur $[t,+\infty[$. Elle change de convexité en $t$ et $f(t)=\frac{3}{2t}$ : le point d'inflexion est $I\pt{t;\frac{3}{2t}}$.

\item \textit{[\num{0,5} pt]} Posons $m=\frac{e^{1}}{2}$ et $s=e^{-1}$. La fonction est continue sur $]0,+\infty[$. Sur $]0,s]$, elle est strictement croissante, de limite $-\infty$ en $0$ et de valeur $2m$ en $s$ : le corollaire du TVI donne exactement une solution dans $]0,s[$. Sur $[s,+\infty[$, elle est strictement décroissante de $2m$ vers $0<m$ : exactement une solution dans $]s,+\infty[$. Ces intervalles couvrent le domaine et $f(s)\ne m$ : il y a exactement deux solutions.

\item \textit{[\num{0,5} pt]} Posons $F(x)=2\ln x+\frac12(\ln x)^2$, dérivable pour $x>0$. Alors $F'(x)=\frac{2}{x}+\frac{\ln x}{x}=f(x)$ ; $F$ est une primitive sur $]0,+\infty[$.

\item \textit{[\num{0,5} pt]} Avec cette primitive, $F(e)=\frac{5}{2}$ et $F(1)=0$. L'intégrale vaut donc $\frac{5}{2}$, et la valeur moyenne est
\[\mu=\frac{1}{e-1}\int_1^e f(x)\,\dd x=\frac{5}{2(e-1)}.\]

\item \textit{[1 pt]} Domaine $x>2$ ; $x^2-4=21$, donc $x=5$, la racine négative étant exclue.

\end{enumerate}

\subsection*{Exercice 3 · Droite, plan et projeté}\begin{enumerate}

\item \textit{[\num{0,5} pt]} $\pt{x;y;z}=\pt{2+t;2t;2t}$, $t\in\R$.

\item \textit{[\num{0,5} pt]} $2+9t=10$ donne $t=\frac{8}{9}$ ; $H=\pt{\frac{26}{9};\frac{16}{9};\frac{16}{9}}$.

\item \textit{[\num{0,5} pt]} Le vecteur $\vec n=\pt{1;2;2}$ est normal à $\mathcal P$. Comme $H\in\mathcal P$ et $\vect{MH}=\frac{8}{9}\vec n$, $H$ est le projeté orthogonal de $M$ sur le plan. Le repère étant orthonormé, $\|\vec n\|=\sqrt{1+4+4}=3$. Donc
\[d(M,\mathcal P)=MH=|\frac{8}{9}|\,3=\frac{8}{3}.\]

\item \textit{[1 pt]} Poser $y=z=s$ ; $x=10-4s$. Donc $\pt{x;y;z}=\pt{10;0;0}+s\pt{-4;1;1}$.

\item \textit{[\num{0,5} pt]} Oui, ses deux dernières coordonnées sont égales. $H$ est dans la droite et $\vect{MH}$ est orthogonal à toute direction du plan, donc à cette droite. Distance $\frac{8}{3}$.

\item \textit{[\num{0,5} pt]} $x+2y+2z=2$ ; même vecteur normal et passage par $M$.

\item \textit{[\num{0,5} pt]} $\vec n\cdot\pt{-4;1;1}=-4+2+2=0$ : angle $\frac{\pi}{2}$.

\item \textit{[1 pt]} $N=M+\vec n$ est sur la même droite normale que $M$ ; son projeté sur $\mathcal P$ est donc aussi $H$. Or $H\in\mathcal P\cap\mathcal Q$. Pour tout $P$ de cette droite, $NP^2=NH^2+HP^2$, donc sa distance minimale à $N$ est $NH$. Finalement,
\[d(N,\mathcal P)=d(N,\mathcal P\cap\mathcal Q)=\|\vect{NH}\|=|\frac{8}{9}-1|\,\|\vec n\|=\frac{1}{3}.\]

\end{enumerate}

\subsection*{Exercice 4 · Répétitions et estimation}\begin{enumerate}

\item \textit{[\num{0,5} pt]} Répétition d'un nombre fixe $18$ d'épreuves identiques indépendantes à deux issues ; $X\sim\mathcal B\pt{18;\frac{1}{5}}$.

\item \textit{[\num{0,5} pt]} $P(X=2)=\binom{18}2(\frac{1}{5})^2(\frac{4}{5})^{16}\simeq\num{0,172263}$.

\item \textit{[\num{0,5} pt]} $1-P(X=0)-P(X=1)=1-(\frac{4}{5})^{18}-18(\frac{1}{5})(\frac{4}{5})^{17}\simeq\num{0,900921}$.

\item \textit{[1 pt]} $E(X)=\frac{18}{5}$, $V(X)=\frac{72}{25}$.

\item \textit{[\num{0,5} pt]} Pour $N$ essais indépendants, la probabilité d'aucun succès est $(1-p)^N$. La probabilité demandée est donc $1-(\frac{4}{5})^N$. Elle est strictement croissante avec $N$, car $0<\frac{4}{5}<1$. Les calculs donnent
\[P_{13}\simeq\num{0,945024}<\num{0,95},\qquad P_{14}\simeq\num{0,956020}\ge\num{0,95}.\]
Ainsi $14$ est le plus petit nombre d'essais convenable.

\item \textit{[\num{0,5} pt]} Notons $n=18$. Pour $j\in\{0,\ldots,n\}$, un résultat ordonné contenant $j$ succès et $n-j$ échecs a, par indépendance, la probabilité $p^j(1-p)^{n-j}$. Il existe $\binom nj$ choix des positions des succès. Ces résultats sont incompatibles, donc
\[P(X=j)=\binom nj p^j(1-p)^{n-j},\qquad p=\frac{1}{5}.\]

\item \textit{[\num{0,5} pt]} Les événements $X=1$, $X=2$, $X=3$ sont incompatibles, donc
\[P(1\le X\le3)=\sum_{j=1}^3\binom{18}j(\frac{1}{5})^j(\frac{4}{5})^{18-j}\simeq\num{0,483011}.\]

\mxneed{24}

\item \textit{[1 pt]} 
\begin{lstlisting}
from math import comb

def quantile():
    n, p = 18, 2/10
    j, cumul = 0, (1-p)**n
    while cumul < 0.95:
        j += 1
        cumul += comb(n, j)*p**j*(1-p)**(n-j)
    return j
\end{lstlisting}
Les probabilités ajoutées sont positives et le cumul final vaut $1$, donc la boucle s'arrête. Elle renvoie $7$. Les cumuls aux rangs $6$ et $7$ sont environ $\num{0,948729}$ et $\num{0,983720}$ : cela confirme la minimalité.

\end{enumerate}

\fi

\fi

\ifmxteacher

\mxnewpage\phantomsection\label{sujet:8-4}\addcontentsline{toc}{chapter}{Blanc 2 · B}

\section*{Blanc 2 · B}

\mxbanner*{SUJET RÉSERVÉ À L'ÉVALUATION}

\textbf{Durée : 4 h · Note sur 20 · \mxTallied{ex-8-4} exercices indépendants.}\\Calculatrice autorisée. Toute réponse doit être justifiée ; les démarches incomplètes seront prises en compte. Donner les valeurs exactes ; arrondir les probabilités à $10^{-6}$ et les autres approximations à $10^{-4}$, sauf précision différente.

\smallskip\textbf{Nom :}\enspace\dotfill\quad\textbf{Classe :}\enspace\dotfill\quad\textbf{Date :}\enspace\dotfill

\mxtally{ex-8-4}
\subsection*{Exercice 1 · Intégrale et suite \hfill 5 points}Pour tout entier naturel $n$, on pose $I_n=\displaystyle\int_0^1x^ne^{2x}\,\dd x$. Pour $n=0$, l'intégrande est la fonction $x\mapsto e^{2x}$.

\begin{enumerate}

\item Calculer $I_0$. \hfill\textit{(\num{0,5} pt)}

\item À l'aide d'une intégration par parties, calculer $I_1$. \hfill\textit{(\num{0,5} pt)}

\item Démontrer, en comparant les intégrandes, que $(I_n)$ est positive et décroissante. \hfill\textit{(\num{0,5} pt)}

\item Montrer que $I_n\le\frac{e^{2}}{n+1}$ pour tout $n\in\N$. En déduire la limite de $(I_n)$. \hfill\textit{(1 pt)}

\item À l'aide d'une intégration par parties, exprimer $I_{n+1}$ en fonction de $I_n$ et de $n$. \hfill\textit{(\num{0,5} pt)}

\item Montrer que $I_n\ge\frac1{n+1}$ pour tout $n\in\N$. \hfill\textit{(\num{0,5} pt)}

\item On partage $\intff{0}{1}$ en $N$ intervalles de même longueur, avec $N\ge1$. Écrire les sommes $L_N$ et $U_N$ des aires des rectangles construits avec les valeurs de $e^{2x}$ aux extrémités gauches et droites. Justifier $L_N\le I_0\le U_N$, puis calculer $U_N-L_N$. \hfill\textit{(\num{0,5} pt)}

\item Déterminer le plus petit entier $N$ garantissant que cet encadrement a une largeur au plus égale à $10^{-2}$. Écrire une fonction Python qui renvoie $N$, $L_N$ et $U_N$. \hfill\textit{(1 pt)}

\end{enumerate}

\mxtally{ex-8-4}
\subsection*{Exercice 2 · Une fonction logarithmique \hfill 5 points}On considère la fonction $f$ définie sur $\intoo{0}{+\infty}$ par $f(x)=\frac{5+\ln x}x$. On note $\mathcal C_f$ sa courbe représentative.

\begin{enumerate}

\item Déterminer les limites de $f$ en $0^+$ et en $+\infty$. \hfill\textit{(\num{0,5} pt)}

\item Calculer $f' (x)$, puis dresser le tableau de variations de $f$. \hfill\textit{(\num{0,5} pt)}

\item Préciser la valeur du maximum de $f$, puis résoudre $f(x)=0$. \hfill\textit{(\num{0,5} pt)}

\item Étudier la convexité de $f$ et déterminer les coordonnées du point d'inflexion de $\mathcal C_f$. \hfill\textit{(1 pt)}

\item Démontrer que l'équation $f(x)=\frac12 e^{4}$ admet exactement deux solutions dans $\intoo{0}{+\infty}$. \hfill\textit{(\num{0,5} pt)}

\item Déterminer une primitive de $f$ sur $\intoo{0}{+\infty}$. \hfill\textit{(\num{0,5} pt)}

\item Calculer la valeur moyenne de $f$ sur $\intff{1}{e}$. \hfill\textit{(\num{0,5} pt)}

\item Indépendamment des questions précédentes, résoudre $\ln(x-3)+\ln(x+3)=\ln(27)$, après avoir précisé l'ensemble sur lequel cette équation est définie. \hfill\textit{(1 pt)}

\end{enumerate}

\mxnewpage\mxcontinued{Blanc 2 · B · suite}

\mxtally{ex-8-4}
\subsection*{Exercice 3 · Droite, plan et projeté \hfill 5 points}Dans un repère orthonormé de l'espace, on considère le plan $\mathcal P$ d'équation $x+2y+2z=11$ et le point $M\pt{3;0;0}$. La droite $d$ passe par $M$ et admet pour vecteur directeur $\vec n=\pt{1;2;2}$.

\begin{enumerate}

\item Donner une représentation paramétrique de $d$. \hfill\textit{(\num{0,5} pt)}

\item Démontrer que $d$ coupe $\mathcal P$ en un point $H$, puis calculer les coordonnées de $H$. \hfill\textit{(\num{0,5} pt)}

\item Justifier que $H$ est le projeté orthogonal de $M$ sur $\mathcal P$, puis calculer la distance de $M$ à $\mathcal P$. \hfill\textit{(\num{0,5} pt)}

\item On considère le plan $\mathcal Q$ d'équation $y=z$. Démontrer que $\mathcal P\cap\mathcal Q$ est une droite et en donner une représentation paramétrique. \hfill\textit{(1 pt)}

\item Vérifier que $H\in\mathcal Q$. En déduire la distance de $M$ à la droite $\mathcal P\cap\mathcal Q$. \hfill\textit{(\num{0,5} pt)}

\item Déterminer une équation cartésienne du plan parallèle à $\mathcal P$ passant par $M$. \hfill\textit{(\num{0,5} pt)}

\item Déterminer l'angle entre $d$ et la droite $\mathcal P\cap\mathcal Q$. \hfill\textit{(\num{0,5} pt)}

\item On définit le point $N$ par $\vect{MN}=\vec n$. Calculer les distances de $N$ au plan $\mathcal P$ et à la droite $\mathcal P\cap\mathcal Q$, puis les comparer. \hfill\textit{(1 pt)}

\end{enumerate}

\mxtally{ex-8-4}
\subsection*{Exercice 4 · Répétitions et estimation \hfill 5 points}On répète $20$ fois une épreuve à deux issues, appelées succès et échec. Les épreuves sont indépendantes et la probabilité de succès est $p=\frac{1}{10}$ à chaque épreuve. On note $X$ le nombre de succès obtenus.

\begin{enumerate}

\item Justifier que $X$ suit une loi binomiale et préciser ses paramètres. \hfill\textit{(\num{0,5} pt)}

\item Calculer la probabilité d'obtenir exactement deux succès. \hfill\textit{(\num{0,5} pt)}

\item Exprimer à l'aide des événements contraires, puis calculer, la probabilité d'obtenir au moins deux succès. \hfill\textit{(\num{0,5} pt)}

\item Calculer l'espérance et la variance de $X$. \hfill\textit{(1 pt)}

\item On répète maintenant la même épreuve un nombre $N\ge1$ de fois, toujours indépendamment. Déterminer le plus petit entier $N$ pour lequel la probabilité d'obtenir au moins un succès est supérieure ou égale à $\num{0,95}$. Une recherche par tableau ou par boucle est autorisée. \hfill\textit{(\num{0,5} pt)}

\item Pour un entier $j\in\{0,\ldots,20\}$, démontrer la formule de $P(X=j)$ en dénombrant les positions possibles des succès. \hfill\textit{(\num{0,5} pt)}

\item On accepte une série de $20$ épreuves si son nombre de succès est compris entre $1$ et $3$, bornes incluses. Calculer la probabilité qu'une série soit acceptée. \hfill\textit{(\num{0,5} pt)}

\item Écrire une fonction Python sans argument qui renvoie le plus petit entier $j$ tel que $P(X\le j)\ge\num{0,95}$. Exécuter la fonction et vérifier la minimalité du résultat. \hfill\textit{(1 pt)}

\end{enumerate}

\ifmxteacher\mxnewpage

\section*{Corrigé · Blanc 2 · B}\label{cor:8-4}

\subsection*{Exercice 1 · Intégrale et suite}\begin{enumerate}

\item \textit{[\num{0,5} pt]} La fonction $x\mapsto \frac{e^{2x}}{2}$ est une primitive de $e^{2x}$. Ainsi $I_0=[\frac{e^{2x}}{2}]_0^1=\frac{e^{2}-1}{2}$.

\item \textit{[\num{0,5} pt]} Choisissons $u(x)=x$ et $v(x)=\frac{e^{2x}}{2}$, deux fonctions dérivables à dérivées continues sur $[0,1]$. Alors $u'=1$ et $v'=e^{2x}$. L'intégration par parties donne
\[\begin{aligned}I_1&=\left[\frac{xe^{2x}}{2}\right]_0^1-\frac{1}{2}\int_0^1e^{2x}\,\dd x\\&=\frac{e^{2}}{2}-\frac{e^{2}-1}{4}=\frac{e^{2}+1}{4}.\end{aligned}\]

\item \textit{[\num{0,5} pt]} Pour $x\in[0,1]$, $0\le x^{n+1}\le x^n$. Comme $e^{2x}>0$, l'intégration sur des bornes croissantes donne $0\le I_{n+1}\le I_n$. La suite est donc positive et décroissante au sens large. Les inégalités sont même strictes : les intégrandes de $I_n$ et de $I_n-I_{n+1}$ sont continues et strictement positives sur $]0,1[$.

\item \textit{[1 pt]} $e^{2x}\le e^{2}$ ; donc $0\le I_n\le e^{2}\int_0^1x^n\,\dd x=\frac{e^{2}}{n+1}\to0$.

\item \textit{[\num{0,5} pt]} Appliquons l'IPP à $u(x)=x^{n+1}$ et $v(x)=\frac{e^{2x}}{2}$. Comme $n+1\ge1$, le terme en $x=0$ est nul. D'où
\[\begin{aligned}I_{n+1}&=\left[\frac{x^{n+1}e^{2x}}{2}\right]_0^1-\frac{n+1}{2}\int_0^1x^ne^{2x}\,\dd x\\&=\frac{e^{2}}{2}-\frac{n+1}{2}I_n.\end{aligned}\]

\item \textit{[\num{0,5} pt]} Pour $0\le x\le1$, $e^{2x}\ge1$ et $x^n\ge0$. Donc $x^n\le x^ne^{2x}$ ; par intégration,
\[I_n\ge\int_0^1x^n\,\dd x=\left[\frac{x^{n+1}}{n+1}\right]_0^1=\frac{1}{n+1}.\]

\item \textit{[\num{0,5} pt]} La fonction $x\mapsto e^{2x}$ est croissante. Sur chaque intervalle $[\frac{j}{N},\frac{j+1}{N}]$, les valeurs aux extrémités encadrent la fonction. En intégrant puis en sommant,
\[L_N=\frac{1}{N}\sum_{j=0}^{N-1}e^{\frac{2j}{N}}\le I_0\le U_N=\frac{1}{N}\sum_{j=1}^Ne^{\frac{2j}{N}}.\]
Les termes intérieurs s'annulent dans la différence, donc $U_N-L_N=\frac{e^{2}-1}{N}$.

\mxneed{24}

\item \textit{[1 pt]} La largeur est au plus $10^{-2}$ si $N\ge100(e^{2}-1)$. Le premier entier satisfaisant cette borne est $N=639$. Voici le calcul des bornes :
\begin{lstlisting}
from math import exp, ceil

def bornes():
    N = ceil(100*(exp(2)-1))
    bas = sum(exp(2*j/N) for j in range(N))/N
    haut = bas+(exp(2)-1)/N
    return N, bas, haut
\end{lstlisting}


\end{enumerate}

\subsection*{Exercice 2 · Une fonction logarithmique}\begin{enumerate}

\item \textit{[\num{0,5} pt]} En $0^+$, le numérateur tend vers $-\infty$ et le dénominateur vers $0^+$ : limite $-\infty$ ; en $+\infty$, limite $0$ par croissance comparée.

\item \textit{[\num{0,5} pt]} Par la règle du quotient,
\[f'(x)=\frac{(\frac{1}{x})x-(5+\ln x)}{x^2}=\frac{-4-\ln x}{x^2}.\]
Le dénominateur est positif. Le numérateur s'annule en $s=e^{-4}$, est positif pour $0<x<s$ et négatif pour $x>s$. Ainsi $f$ croît strictement sur $]0,s]$, puis décroît strictement sur $[s,+\infty[$.

\item \textit{[\num{0,5} pt]} Le maximum est atteint uniquement en $s=e^{-4}$ et vaut
\[f(s)=\frac{5- 4}{e^{-4}}=e^{4}.\]
Comme $x>0$, $f(x)=0$ équivaut à $5+\ln x=0$, soit $x=e^{-5}$.

\item \textit{[1 pt]} On dérive $(-4-\ln x)x^{-2}$ :
\[f''(x)=-x^{-3}-2(-4-\ln x)x^{-3}=\frac{2\ln x+7}{x^3}.\]
Le numérateur s'annule en $t=e^{-\frac{7}{2}}$ et passe du signe négatif au signe positif. La fonction est concave sur $]0,t]$, convexe sur $[t,+\infty[$. Elle change de convexité en $t$ et $f(t)=\frac{3}{2t}$ : le point d'inflexion est $I\pt{t;\frac{3}{2t}}$.

\item \textit{[\num{0,5} pt]} Posons $m=\frac{e^{4}}{2}$ et $s=e^{-4}$. La fonction est continue sur $]0,+\infty[$. Sur $]0,s]$, elle est strictement croissante, de limite $-\infty$ en $0$ et de valeur $2m$ en $s$ : le corollaire du TVI donne exactement une solution dans $]0,s[$. Sur $[s,+\infty[$, elle est strictement décroissante de $2m$ vers $0<m$ : exactement une solution dans $]s,+\infty[$. Ces intervalles couvrent le domaine et $f(s)\ne m$ : il y a exactement deux solutions.

\item \textit{[\num{0,5} pt]} Posons $F(x)=5\ln x+\frac12(\ln x)^2$, dérivable pour $x>0$. Alors $F'(x)=\frac{5}{x}+\frac{\ln x}{x}=f(x)$ ; $F$ est une primitive sur $]0,+\infty[$.

\item \textit{[\num{0,5} pt]} Avec cette primitive, $F(e)=\frac{11}{2}$ et $F(1)=0$. L'intégrale vaut donc $\frac{11}{2}$, et la valeur moyenne est
\[\mu=\frac{1}{e-1}\int_1^e f(x)\,\dd x=\frac{11}{2(e-1)}.\]

\item \textit{[1 pt]} Domaine $x>3$ ; $x^2-9=27$, donc $x=6$, la racine négative étant exclue.

\end{enumerate}

\subsection*{Exercice 3 · Droite, plan et projeté}\begin{enumerate}

\item \textit{[\num{0,5} pt]} $\pt{x;y;z}=\pt{3+t;2t;2t}$, $t\in\R$.

\item \textit{[\num{0,5} pt]} $3+9t=11$ donne $t=\frac{8}{9}$ ; $H=\pt{\frac{35}{9};\frac{16}{9};\frac{16}{9}}$.

\item \textit{[\num{0,5} pt]} Le vecteur $\vec n=\pt{1;2;2}$ est normal à $\mathcal P$. Comme $H\in\mathcal P$ et $\vect{MH}=\frac{8}{9}\vec n$, $H$ est le projeté orthogonal de $M$ sur le plan. Le repère étant orthonormé, $\|\vec n\|=\sqrt{1+4+4}=3$. Donc
\[d(M,\mathcal P)=MH=|\frac{8}{9}|\,3=\frac{8}{3}.\]

\item \textit{[1 pt]} Poser $y=z=s$ ; $x=11-4s$. Donc $\pt{x;y;z}=\pt{11;0;0}+s\pt{-4;1;1}$.

\item \textit{[\num{0,5} pt]} Oui, ses deux dernières coordonnées sont égales. $H$ est dans la droite et $\vect{MH}$ est orthogonal à toute direction du plan, donc à cette droite. Distance $\frac{8}{3}$.

\item \textit{[\num{0,5} pt]} $x+2y+2z=3$ ; même vecteur normal et passage par $M$.

\item \textit{[\num{0,5} pt]} $\vec n\cdot\pt{-4;1;1}=-4+2+2=0$ : angle $\frac{\pi}{2}$.

\item \textit{[1 pt]} $N=M+\vec n$ est sur la même droite normale que $M$ ; son projeté sur $\mathcal P$ est donc aussi $H$. Or $H\in\mathcal P\cap\mathcal Q$. Pour tout $P$ de cette droite, $NP^2=NH^2+HP^2$, donc sa distance minimale à $N$ est $NH$. Finalement,
\[d(N,\mathcal P)=d(N,\mathcal P\cap\mathcal Q)=\|\vect{NH}\|=|\frac{8}{9}-1|\,\|\vec n\|=\frac{1}{3}.\]

\end{enumerate}

\subsection*{Exercice 4 · Répétitions et estimation}\begin{enumerate}

\item \textit{[\num{0,5} pt]} Répétition d'un nombre fixe $20$ d'épreuves identiques indépendantes à deux issues ; $X\sim\mathcal B\pt{20;\frac{1}{10}}$.

\item \textit{[\num{0,5} pt]} $P(X=2)=\binom{20}2(\frac{1}{10})^2(\frac{9}{10})^{18}\simeq\num{0,285180}$.

\item \textit{[\num{0,5} pt]} $1-P(X=0)-P(X=1)=1-(\frac{9}{10})^{20}-20(\frac{1}{10})(\frac{9}{10})^{19}\simeq\num{0,608253}$.

\item \textit{[1 pt]} $E(X)=2$, $V(X)=\frac{9}{5}$.

\item \textit{[\num{0,5} pt]} Pour $N$ essais indépendants, la probabilité d'aucun succès est $(1-p)^N$. La probabilité demandée est donc $1-(\frac{9}{10})^N$. Elle est strictement croissante avec $N$, car $0<\frac{9}{10}<1$. Les calculs donnent
\[P_{28}\simeq\num{0,947665}<\num{0,95},\qquad P_{29}\simeq\num{0,952899}\ge\num{0,95}.\]
Ainsi $29$ est le plus petit nombre d'essais convenable.

\item \textit{[\num{0,5} pt]} Notons $n=20$. Pour $j\in\{0,\ldots,n\}$, un résultat ordonné contenant $j$ succès et $n-j$ échecs a, par indépendance, la probabilité $p^j(1-p)^{n-j}$. Il existe $\binom nj$ choix des positions des succès. Ces résultats sont incompatibles, donc
\[P(X=j)=\binom nj p^j(1-p)^{n-j},\qquad p=\frac{1}{10}.\]

\item \textit{[\num{0,5} pt]} Les événements $X=1$, $X=2$, $X=3$ sont incompatibles, donc
\[P(1\le X\le3)=\sum_{j=1}^3\binom{20}j(\frac{1}{10})^j(\frac{9}{10})^{20-j}\simeq\num{0,745470}.\]

\mxneed{24}

\item \textit{[1 pt]} 
\begin{lstlisting}
from math import comb

def quantile():
    n, p = 20, 1/10
    j, cumul = 0, (1-p)**n
    while cumul < 0.95:
        j += 1
        cumul += comb(n, j)*p**j*(1-p)**(n-j)
    return j
\end{lstlisting}
Les probabilités ajoutées sont positives et le cumul final vaut $1$, donc la boucle s'arrête. Elle renvoie $4$. Les cumuls aux rangs $3$ et $4$ sont environ $\num{0,867047}$ et $\num{0,956826}$ : cela confirme la minimalité.

\end{enumerate}

\fi

\fi

